\clearpage
\fancyhead[L]{\small LOGARITHMENTAFELN.}\pdfbookmark[0]{Bemerkungen über Logarithmentafeln}{logarithmentafeln}
\begin{center}
BEMERKUNGEN ÜBER LOGARITHMENTAFELN.
\end{center}

\begin{center}
\rule{0.65\linewidth}{0.3pt}\\
Monatliche Correspondenz, herausg. vom Freih. v. \textsc{Zach}. 1802 Nov.\\
\rule{0.65\linewidth}{0.3pt}
\end{center}

\begin{center}2.\end{center}

Zu dem S. 497 des vorigen Heftes gegebenen Verzeichniss aller Druckfehler der Stereotype-Ausgabe der \textsc{Callet}'schen logarithm. Tafeln hat Dr. \textsc{Gauss} die
Güte gehabt, noch folgende Errata anzuzeigen:

\begin{center}
\footnotesize\setlength{\tabcolsep}{3pt}\renewcommand{\arraystretch}{1.4}\begin{tabular}{@{}p{0.30\linewidth}lll@{}}
\textit{Log. Sin. de Seconde en Seconde}
& \(4^{0}\;15'\;5''\quad\sin 8.8690096\) & lies & \(8.8700096\) \\
& \(4\quad15\quad6\quad\sin 8.8690379\) & --- & \(8.8700379\) \\
\textit{Log. Sin. de} \(10\) \textit{en} \(10\) \textit{Secondes}
& \(\mathrm{Arc}\,21^{0}\;27'\;20''\) & lies & \(21^{0}\;27'\;30''\).
\end{tabular}
\end{center}

Für \(33^{0}\) unten statt 59 Deg. lies 56 Deg. (nur in einigen Abdrücken).

\Needspace{9\baselineskip}\pdfbookmark[1]{Göttingische gelehrte Anzeigen. 1811 Mai 25.}{prasse-1811}
\begin{center}
\rule{0.60\linewidth}{0.3pt}\\
Göttingische gelehrte Anzeigen. 1811 Mai 25.\\
\rule{0.60\linewidth}{0.3pt}
\end{center}

\textit{Logarithmische Tafeln für die Zahlen, Sinus und Tangenten, neu geordnet
von \textup{\textsc{Moritz von Prasse}}, ordentlichem Professor der Mathematik zu Leipzig.} 80 Seiten in Octav. Leipzig. In Commission bei P. J. Besson.

Diese Tafeln enthalten dasselbe. was die beliebten kleinen Tafeln von \textsc{Lalande} haben, nemlich die Logarithmen aller Zahlen bis 10000, und die Logarith\SourceLeaf{251}%
men der Sinus und Tangenten für alle einzelnen Minuten des Quadranten; alles
auf fünf Decimalen. Allein Hr. v. \textsc{Prasse} hat dieses bei einem nicht viel grössern
Format auf den dritten Theil der Seitenzahl reducirt, indem er die bei den grössern
Tafeln übliche Einrichtung anwandte, immer je zehn Logarithmen in Eine Zeile,
und die ersten Ziffern nur Einmal anzusetzen, wobei aber alle Differenzen haben
wegbleiben müssen. Es scheint also hierdurch an Bequemlichkeit wieder verloren zu gehen, was an Kürze gewonnen wird. Da indessen hierüber nur nach
wirklichem Gebrauche geurtheilt werden kann, so hat Rec., der sich an die kleinen \textsc{Lalande}'schen Tafeln gewöhnt hat, diese eine Zeitlang bei Seite gelegt, und
sich der vorliegenden zu bedienen versucht. Er hat gefunden, dass jene kleinen
Unbequemlichkeiten von dem Vortheile, viel weniger blättern zu müssen, bei den
Logarithmen der Zahlen, die hier auf 31, bei \textsc{Lalande} auf 111 Seiten stehen,
merklich überwogen werden, und er bedient sich daher dieser neuen Tafeln gern.
Nicht so hat er es bei den trigonometrischen Tafeln gefunden, die hier 40, bei
\textsc{Lalande} 90 Seiten einnehmen, besonders desswegen, weil bei Hrn. v. \textsc{Prasse} die
Bogen von 0 bis 90 Grad fortlaufen, und daher die Sinus und Tangenten von den
Cosinus und Cotangenten getrennt sind. Diess ist um so beschwerlicher, da die
Fälle so sehr häufig sind, wo man z. B. von einem Bogen den Sinus und Cosinus
zugleich nöthig hat, oder wo man, ohne den Bogen selbst zu brauchen, aus dem
Sinus oder der Tangente den Cosinus verlangt. Hier würde er also allemal die
\textsc{Lalande}'schen Tafeln vorziehen, und er hätte gewünscht, dass Hr. v. \textsc{Prasse} lieber jede Seite noch einmal in der Mitte durch eine Horizontallinie getheilt hätte,
um jene unangenehme Trennung zu vermeiden, wobei die Zusammendrängung in
den kleinen Raum doch hätte Statt finden können.

Ausserdem unterscheiden sich diese Tafeln noch dadurch, dass allemal die
letzte Ziffer eines jeden Logarithmen, wenn sie vergrössert worden ist, mit einer
andern Schrift gesetzt ist. Hr. v. P. glaubt dadurch grössere Genauigkeit bei den
Rechnungen befördern zu können. Allein da man doch meistens in der Ausübung
nur mit Logarithmen zu rechnen hat, die interpolirt werden müssen, so kann
man nicht ohne Beschwerde auf jenen Umstand Rücksicht nehmen, und so oft
man glaubt, dass die nur auf eine halbe Einheit in der fünften Decimale zuverlässigen Logarithmen nicht genug genaue Resultate geben können, so thut man
besser, grössere Tafeln mit sechs oder sieben Decimalen anzuwenden. Rec. kann
daher diese Einrichtung, die, allgemein zu reden, allerdings die Genauigkeit der
\SourceLeaf{252}
Rechnung zu \emph{verdoppeln} dienen kann, nicht für sehr nützlich anerkennen, zumal
da die Cursivzahlen neben den andern dem Auge unangenehm, und hin und wieder nicht scharf genug sind. Sonst ist der Druck nett; nur werden diejenigen,
die dergleichen Tafeln viel brauchen, stärkeres Papier wünschen.

\par\Needspace{10\baselineskip}\pdfbookmark[1]{Prasse. 1814 December 19.}{prasse1814}
\begin{center}
\rule{0.63\linewidth}{0.3pt}\\
Göttingische gelehrte Anzeigen. 1814 December 19.\\
\rule{0.63\linewidth}{0.3pt}
\end{center}

\emph{Tables logarithmiques pour les nombres, les sinus et les tangentes, disposées
dans un nouvel ordre par M. \textup{\textsc{de Prasse}}, professeur des mathématiques à Berlin} (zu
Leipzig) \emph{corrigées et précédées d’une introduction traduite de l’allemand et accompagnée de notes et d’un avertissement par} M. \textsc{Halma} 1814. Paris. 80 Seiten. Preis
Ein Frank.

Ein neuer Abdruck der geschmeidigen von \textsc{Prasse}’schen Tafeln, welche wir
in diesen Blättern [1811 Mai 25] angezeigt haben. Unser dortiges Urtheil über
die von dem französischen Herausgeber unverändert beibehaltene Anordnung der
Tafeln haben wir durch einen drei Jahre länger fortgesetzten Gebrauch derselben
in allen Stücken bestätigt gefunden. In Ansehung der Schönheit des Drucks und
Papiers scheint uns die Französische Ausgabe der Deutschen eher nachzustehen;
doch sind mehrere Druckfehler der letztern hier berichtigt. Wenn man übrigens
bei einem Werke dieser Art, das der verstorbene von \textsc{Prasse} gewiss nicht Gewinnes halber, sondern zum Dienste der Wissenschaft auf seine Kosten unternahm,
auch nicht weiter untersuchen will, in wie fern Hr. \textsc{Halma} zu einem neuen Abdrucke berechtigt war, so kann man doch nicht umhin, sich zu wundern, dass
derselbe, aus Besorgniss seinerseits wieder nachgedruckt zu werden, die einzelnen Exemplare mit seinem Namenszuge bezeichnet hat, und die Nachdrucker gerichtlich zu belangen droht.
\SourceLeaf{253}
\par\Needspace{10\baselineskip}\pdfbookmark[1]{Tafel zur bequemern Berechnung. 1812.}{gauss1812log}
\begin{center}
\rule{0.66\linewidth}{0.3pt}\\
Monatliche Correspondenz, herausg. vom Freih. v. \textsc{Zach}. 1812. Nov.\\
\rule{0.66\linewidth}{0.3pt}
\end{center}

\emph{Tafel zur bequemern Berechnung des Logarithmen der Summe oder Differenz
zweier Grössen, welche selbst nur durch ihre Logarithmen gegeben sind. Von Herrn
Prof.} \textsc{Gauss}.

Je weiter sich beständig die Geschäfte der rechnenden Astronomen ausdehnen, desto wichtiger wird ihnen jede, wenn auch an sich nur kleine Erleichterung derselben. Die \emph{Monatliche Correspondenz} hat sich hierin schon vielfältige
Verdienste erworben, indem sie mancherlei Tafeln aufgenommen hat, deren kleiner Umfang nicht verstattete, sie besonders herauszugeben. Ich lege daher gern
in derselben eine kleine Tafel nieder, die freilich nicht eigentlich astronomisch
ist, aber besonders doch den rechnenden Astronomen willkommen sein wird, und
die etwa in Zukunft sehr zweckmässig mit einem neuen Abdruck der kleinen \textsc{La
Lande}’schen Tafeln verbunden werden könnte. Das Geschäft, was sie erleichtern soll, kommt bei astronomischen Rechnungen alle Augenblick vor; es erfordert sonst ein dreimaliges, oder wenn man eine leichte Verwandlung anwendet,
doch nothwendig ein zweimaliges Aufschlagen in den Logarithmen-Tafeln, was
hier auf ein einziges gebracht wird. Die Idee dazu hat \textsc{Leonelli}, so viel ich weiss,
zuerst angegeben; allein seine Meinung war, eine solche Tafel für Rechnungen
mit 14 Decimalen zu construiren, und gerade dies kann ich nicht zweckmässig
finden. Sie würde bei einer solchen Ausdehnung einen grossen Folioband füllen,
ihre Berechnung würde eine ungeheuere Arbeit und Zeit erfordern, und sie würde
fast nie von Nutzen, und immer nur von wenig Nutzen sein, da so scharfe Rechnungen so selten --- in der eigentlichen practischen Astronomie nie --- vorkommen, dass die verhältnissmässig doch nur kleine Erleichterung die Construction,
ja nicht einmal den Ankauf einer solchen Tafel belohnen würde. Ich habe diese
Tafel zu meinem eigenen Gebrauch für Rechnungen mit 5 Decimalen, die in der
Ausübung die häufigsten sind, schon vor vielen Jahren construirt, und die, wenn
auch \emph{jedesmal} kleine, doch wenn sie viele Tausendmale wiederkehrt, sehr erhebliche Erleichterung, hat mir die darauf gewandte Mühe bereits reichlich ersetzt.
Es wäre zu wünschen, dass jemand sich der Arbeit unterzöge, eine ähnliche Tafel in 10 oder 100 mal so grosser Ausdehnung für Rechnungen mit 7 Decimalen zu
\SourceLeaf{254}
construiren, die als ein sehr schätzbares Supplement den gewöhnlichen Logarithmen-Tafeln beigefügt werden könnte.

Die Einrichtung der aus drei Columnen bestehenden Tafel ist sehr einfach.
Die erste Columne \(A\) geht von 0 bis 2 durch alle Tausendtheile, von da bis \(3,4\)
durch alle Hunderttheile, und von \(3,4\) bis \(5,0\) durch alle Zehntheile; mit \(5,0\)
kann die Tafel für 5 Decimalen als geschlossen angesehen werden, da die zweite
Columne für diesen und für grössere Werthe von \(A\) verschwindet, und die Zahlen der dritten Columne denen der ersten gleich werden. Setzt man eine Zahl
der ersten Columne \(A=\log m\), so ist in der zweiten Columne \(B=\log(1+\frac{1}{m})\)
und in der dritten Columne \(C=\log(1+m)\), so dass immer \(C=A+B\). Man
kann also auch die Zahlen der drei Columnen als die doppelten Logarithmen der
Tangenten, Cosecanten und Secanten der Winkel von \(45^{0}\) bis \(90^{0}\) betrachten.
Die Anwendung davon ist nun folgende:

\begin{center}
I.\quad \emph{Aus den Logarithmen zweier Grössen \(a,b\) den Logarithmen der Summe
zu finden.}
\end{center}

Es sei \(\log a\) der grössere Logarithm., man gehe mit \(\log a-\log b\) in die
Columne \(A\) ein, und nehme daneben entweder aus der zweiten Columne \(B\), oder
aus der dritten Columne \(C\). Man hat dann
\[
\begin{aligned}
&\log(a+b)=\log a+B\\
\text{oder}\qquad &\log(a+b)=\log b+C.
\end{aligned}
\]

\begin{center}
II.\quad \emph{Aus den Logarithmen zweier Grössen \(a,b\) den Logarithmen der Differenz
zu finden.}
\end{center}

\emph{Erstens,} ist die Differenz der Logarithmen \(\log a-\log b\) grösser als \(0,30103\),
so suche man dieselbe in \(C\), wodurch man hat
\[
\begin{aligned}
&\log(a-b)=\log a-B\\
\text{oder}\qquad &\log(a-b)=\log b+A
\end{aligned}
\]

\emph{Zweitens,} ist \(\log a-\log b\) kleiner als \(0,30103\), so suche man es in \(B\), wodurch wird
\[
\begin{aligned}
&\log(a-b)=\log a-C\\
\text{oder}\qquad &\log(a-b)=\log b-A
\end{aligned}
\]
\SourceLeaf{255}
Es gibt daher bei jeder Aufgabe zwei Auflösungsarten; man thut aber wohl,
sich an eine bestimmte zu gewöhnen, um sich den Gebrauch der Tafel desto leichter mechanisch zu machen. Mir ist dies bei der jedesmal zuerst angesetzten Manier am bequemsten gefallen.

\Needspace{12\baselineskip}\begin{center}\textit{Beispiele:}\end{center}

I. Aus $\log a=0{,}36173$ und $\log b=0{,}23045$ den Logarithmen der Summe
zu finden, sucht man $0{,}13128$ in $A$, wobei man findet
\[
\begin{array}{rcr}
B&\cdot\;\cdot\;\cdot\;\cdot\;\cdot&0{,}24033\\
\log a&\cdot\;\cdot\;\cdot\;\cdot\;\cdot&0{,}36173\\\hline
\log(a+b)&\cdot\;\cdot&0{,}60206
\end{array}
\qquad
\begin{array}{rcr}
C&\cdot\;\cdot\;\cdot\;\cdot\;\cdot&0{,}37161\\
\log b&\cdot\;\cdot\;\cdot\;\cdot\;\cdot&0{,}23045\\\hline
&&0{,}60206
\end{array}
\]

II. Aus $\log a=0{,}89042$, und $\log b=0{,}24797$ den Logarithm. der Differenz zu finden. Da $\log a-\log b=0{,}64245$ grösser als $0{,}30103$, so sucht
man es in der Columne $C$, woneben man findet
\[
\begin{array}{rcr}
B&\cdot\;\cdot\;\cdot\;\cdot\;\cdot&0{,}11227\\
\log a&\cdot\;\cdot\;\cdot\;\cdot\;\cdot&0{,}89042\\\hline
\log(a-b)&\cdot\;\cdot&0{,}77815
\end{array}
\qquad
\begin{array}{rcr}
A&\cdot\;\cdot\;\cdot\;\cdot\;\cdot&0{,}53018\\
\log b&\cdot\;\cdot\;\cdot\;\cdot\;\cdot&0{,}24797\\\hline
&&0{,}77815
\end{array}
\]

III. Aus $\log a=0{,}25042$, $\log b=0{,}19033$ den Logarithmen der Differenz zu finden. Hier gibt $\log a-\log b=0{,}06009$ in $B$ aufgesucht.
\[
\begin{array}{rcr}
C&\cdot\;\cdot\;\cdot\;\cdot\;\cdot&0{,}88871\\
\log a&\cdot\;\cdot\;\cdot\;\cdot\;\cdot&0{,}25042\\\hline
\log(a-b)&\cdot\;\cdot&9{,}36171
\end{array}
\qquad
\begin{array}{rcr}
A&\cdot\;\cdot\;\cdot\;\cdot\;\cdot&0{,}82862\\
\log b&\cdot\;\cdot\;\cdot\;\cdot\;\cdot&0{,}19033\\\hline
&&9{,}36171
\end{array}
\]

\par\Needspace{10\baselineskip}\pdfbookmark[1]{Pasquich. 1817 October 4.}{pasquich1817}
\begin{center}
\rule{0.65\linewidth}{0.4pt}\\
Göttingische gelehrte Anzeigen. 1817 October 4.\\
\rule{0.65\linewidth}{0.4pt}
\end{center}

\textit{Abgekürzte Logarithmisch-Trigonometrische Tafeln, mit neuen Zusätzen zur
Abkürzung und Erleichterung trigonometrischer Rechnungen, herausgegeben von }\textsc{Joh.}
\textsc{Pasquich}\textit{, Director der Königl. Ofner Sternwarte.} Leipzig. In der Weidmannschen
Buchhandlung 1817. XXXII und 228 Seiten in Octav. (Auch mit Lateinischem Titel.)
\par\SourceLeaf{256}
Kleinere logarithmische Tafeln, mit fünf Decimalen, sind bei denjenigen,
die viel mit Zahlenrechnungen zu verkehren haben, besonders bei den Astronomen, sehr beliebt, weil in der That die Fälle, wo sie ausreichen, häufig. ja die häufigeren, sind, und durch ein bequemes Format und eine mässige Grösse die Arbeit
sehr erleichtert wird. Die kleinen netten \textsc{Lalande}'schen und die \textsc{Prasse}'schen
Tafeln sind in Jedermanns Händen; bei letztern ist das Zusammendrängen in einen kleinen Raum so weit wie möglich getrieben, zum Theil aber allerdings auf
Kosten der Bequemlichkeit. Die Herausgabe der vorliegenden auch nur auf fünf
Stellen gehenden Tafeln ist, wie in der Vorrede berichtet wird, durch den von
\textsc{Gauss} in der monatlichen Correspondenz 1812 geäusserten Wunsch veranlasst,
dass die daselbst zuerst abgedruckte Tafel zur bequemen Berechnung der Logarithmen der Summen und Differenzen einer neuen Ausgabe der \textsc{Lalande}'schen Tafeln einverleibt werden möchte. Der neue Abdruck dieser Hülfstafel in gegenwärtiger Sammlung wird denjenigen angenehm sein, denen der erste Abdruck
nicht zu Gebote stand, oder denen der Gebrauch derselben in der M. C. zu beschwerlich war. Ausserdem zeichnet sich diese Sammlung noch durch eine neue
von Hrn. \textsc{Pasquich} berechnete den trigonometrischen Tafeln beigefügte Hülfstafel aus, deren wir unten mit mehrern erwähnen werden.

Die Logarithmen der Zahlen gehen, wie bei \textsc{Lalande} und von \textsc{Prasse} bis
10000, und sind, so wie bei jenen, hier in ihrer natürlichen Ordnung gedruckt.
Doch vermisst man ungern ein Paar Erleichterungsmittel, welche bei den \textsc{Lalande}'schen Tafeln Statt finden; es sind nemlich theils die \textit{Differenzen} nicht beigefügt, theils die untersten Logarithmen jeder Spalte oben in der nächstfolgenden
nicht wiederholt. In den \textsc{Prasse}'schen Tafeln findet man zwar diese Bequemlichkeit auch nicht, allein dort werden sie durch den kleinen Raum der Tafel mehr
als ersetzt, da jene auf 24 Seiten eben dasselbe liefern, was bei \textsc{Pasquich} auf 56
Seiten eines beträchtlich grössern Formats steht. Dies scheinen zwar nur Kleinigkeiten, und sie sind es auch für alle, die nur dann und wann einmal Logarithmen aufzuschlagen haben, aber nicht für solche, die Logarithmen-Tafeln beständig zur Hand haben müssen.

In den trigonometrischen Tafeln enthält immer jede Seite zur linken die
Logarithmen der Sinus, Cosinus, Tangenten und Cotangenten, und zwar so,
dass je drei Seiten zwei Grade fassen. Diese Einrichtung, welche durch das gewählte Format und die Schrift herbeigeführt wurde, scheint uns etwas unbequem;
\SourceLeaf{257}
wir hätten entweder ein kleineres Format, immer mit einem halben Grad auf der
Seite, oder ein etwas weniges längeres mit kleinerer Schrift, so dass ein ganzer
Grad auf die Seite gekommen wäre (wie in \textsc{Shervins} Tafeln) vorgezogen. Von diesen Logarithmen sind immer nur die vier, drei oder zwei letzten Ziffern, so lange
die vorhergehenden ungeändert bleiben, abgedruckt, wodurch dem Copiisten, dem
Setzer und dem Corrector die Arbeit erleichtert wurde, und die Tafeln ein reinlicheres Ansehen erhalten: dem ungeachtet können wir diese Einrichtung bei
Tafeln, die zum täglichen Gebrauch bestimmt sind, nicht unbedingt billigen, da
das Auge immer die, wenn auch nur kleine, Mühe hat, in der Columne erst in
die Höhe zu gehen, und die übrigen Ziffern zu finden. Die Differenzen der Logarithmen findet man hier sogleich mit 60 dividirt; eine Einrichtung, welche auch
in einigen andern Tafeln gewählt ist, in der Absicht, das Interpoliren zu erleichtern. Ob diese Erleichterung wirklich Statt findet, oder nicht, wird von der Gewöhnung des Rechners abhängen. Rec. findet in dieser Beziehung die \textsc{Lalande}schen Tafeln, wo die ganzen Differenzen angesetzt sind, wenigstens nicht unbequemer. Bei der Kleinheit der Zahlen, mit denen zu operiren ist, macht ein
etwas geübter Rechner die zum Behuf des Interpolirens nöthigen Operationen
leicht im Kopfe, und findet fast immer diesen oder jenen Local-Vortheil zu benutzen Gelegenheit. Dabei hat man noch die angenehme Gewissheit, sein Interpolations-Resultat so scharf zu erhalten, als es möglich ist; bei der von Hrn. \textsc{Pasquich} gewählten Einrichtung hingegen ist, allgemein zu reden, der Fehler des
Interpolirens etwas grösser, welches indessen ausführlicher zu entwickeln hier
nicht der Ort ist.

Die Seite zur rechten enthält bei den trigonometrischen Tafeln die Quadrate
der Sinus, Cosinus, Tangenten und Cotangenten, welche zur Erleichterung des
Interpolirens dienen sollen, wenn man aus dem Logarithmen eines Sinus, Cosinus, einer Tangente oder Cotangente den Logarithmen einer der drei andern trigonometrischen Functionen verlangt, ohne den Bogen selbst nöthig zu haben.
Diese Operation kömmt allerdings äusserst häufig vor, und das gewöhnliche Verfahren erfordert beim Interpoliren eine Multiplication und eine Division, wo mit
Hrn. \textsc{Pasquich}'s Hülfstafel eine Multiplication ausreicht. Es ist nemlich, für das
Interpoliren hinreichend genau,
\SourceLeaf{258}
\[
\begin{aligned}
\Delta\log\cos\varphi&=-\operatorname{tang}\varphi^{2}.\Delta\log\sin\varphi\\
\Delta\log\operatorname{tang}\varphi&=-\Delta\log\operatorname{cotang}\varphi=(1+\operatorname{tang}\varphi^{2}).\Delta\log\sin\varphi\\
\Delta\log\sin\varphi&=-\operatorname{cotang}\varphi^{2}.\Delta\log\cos\varphi\\
\Delta\log\operatorname{tang}\varphi&=-\Delta\log\operatorname{cotang}\varphi=-(1+\operatorname{cotang}\varphi^{2}).\Delta\log\cos\varphi\\
\Delta\log\sin\varphi&=\cos\varphi^{2}.\Delta\log\operatorname{tang}\varphi=-\cos\varphi^{2}.\Delta\log\operatorname{cotang}\varphi\\
\Delta\log\cos\varphi&=-\sin\varphi^{2}.\Delta\log\operatorname{tang}\varphi=\sin\varphi^{2}.\Delta\log\operatorname{cotang}\varphi
\end{aligned}
\]
Inzwischen muss Rec. gestehen, dass er dem ungeachtet das gewöhnliche Verfahren zum Interpoliren nicht bloss eben so bequem, sondern sogar bequemer findet. Theils wird es immer erst einige Mühe kosten, sich die obigen sechs Formeln so mechanisch zu machen, dass man sie ohne alles Besinnen oder ohne ein
besonderes Blatt neben sich zu legen, richtig anwendet; theils ist es beschwerlich, den Multiplications-Factor erst auf der andern Seite aufzusuchen, oder vielmehr zusammen zu suchen, da die oben erwähnte Trennung der ersten und letzten Ziffern auch hier beim Abdruck gewählt ist; endlich hat man bei dem gewöhnlichen Verfahren es immer nur mit kleinen Zahlen zu thun, mit denen man
leicht im Kopf rechnet, da hingegen die Quadrate in \textsc{Pasquich}'s Tafeln mit fünf
Decimalen angesetzt sind, die man freilich nicht alle braucht, aber die gerade deswegen, wie jeder erfahrne Rechner weiss, störend sind. Ausserdem können wir
hier nicht unerwähnt lassen, dass das gewöhnliche Verfahren, allgemein zu reden,
\textit{schärfer} ist, als diese künstlichere Interpolation (die Gründe dieser Behauptung,
von der man vielleicht bei einer weniger genauen Prüfung gerade das Gegentheil
glauben könnte, würden für diesen Ort zu weitläuftig sein). Wir begnügen uns
das Gesagte bloss durch ein Beispiel zu erläutern. Soll zu \(\log\cos\varphi=9{,}92478\)
der \(\log\operatorname{tang}\varphi\) gesucht werden, so findet man den Proportionaltheil aus \textsc{Pasquich}'s Tafel durch die Berechnung von \(4\times(1+2{,}4170)=13{,}668\) oder am nächsten \(=14\), also \(\log\operatorname{tang}\varphi=9{,}80850\), während die gewöhnliche Methode den
Proportionaltheil eben so bequem durch die Entwicklung von \(\frac{4\times28}{9}=12\frac49\), am
nächsten \(=12\), und den gesuchten Logarithmen \(=9{,}80848\) gibt. In diesem
Beispiele ist auch das Resultat der gewöhnlichen Methode das schärfere; in andern Fällen kann auch das umgekehrte Verhältniss Statt finden, aber im Durchschnitt wird der Vortheil in dieser Beziehung auf Seiten des gewöhnlichen Verfahrens sein. Uebrigens wollen wir nicht in Abrede stellen, dass dieser Theil
der Tafel, wenn auch das Interpoliren nicht dadurch gewinnt, doch zuweilen für
\SourceLeaf{259}
andere Zwecke angenehm sein könne; allein die Bequemlichkeit logarithmischer
Handtafeln, die man zum täglichen Gebrauch bestimmt, verliert natürlich in demselben Verhältniss, als ihr Umfang vergrössert wird. Wir bemerken noch, dass
in dem ersten Grade die trigonometrischen Logarithmen von \(10\) zu \(10\) Secunden
bis \(56\) Minuten, und in den vier letzten Minuten von \(20\) zu \(20\) Secunden angesetzt sind.

Die \textsc{Gauss}ische Tafel für die Logarithmen der Summen und Differenzen ist
ganz unverändert abgedruckt. Inconsequent scheint es uns aber zu sein, wenn
der Verf. in der Einleitung den Nutzen einer ähnlichen Tafel mit sieben Decimalen in Zweifel zieht. Ist anders eine solche Tafel zweckmässig eingerichtet, so ist
ihr Nutzen bei scharfen Rechnungen gerade eben so gross, als der Nutzen der
hier wieder abgedruckten Tafeln bei Rechnungen mit fünf Decimalen: bei den
kleinern Tafeln, eben so wie bei den grössern, wird der dadurch zu erhaltende
Zeitgewinn natürlich nur solchen Personen fühlbar, die \textit{viel} zu rechnen haben.
Wir haben jetzt bald die Erscheinung einer solchen grössern Tafel, von einer geschickten Hand berechnet, zu erwarten.

\par\Needspace{10\baselineskip}\pdfbookmark[1]{Matthiessen. 1819 Januar 30.}{matthiessen1819}\medskip\hrule\smallskip
\begin{center}
Göttingische gelehrte Anzeigen. 1819 Januar 30.
\end{center}
\smallskip\hrule\medskip

[E. A. \textsc{Matthiessen}.] \textit{Tafel zur bequemern Berechnung des Logarithmen der
Summe oder Differenz zweier Grössen, welche selbst nur durch ihre Logarithmen gegeben sind.} Altona 1818. Bei J. F. Hammerich. Einleitung 33 S. Die Tafeln
212 S. in Quart. (Titel und Einleitung auch in lateinischer Sprache.)

Bei etwas ausgedehnten Rechnungen ist jede Abkürzung schätzbar; auch
solche Hülfsmittel, die bei einer einmaligen Anwendung nur einen kleinen Vortheil gewähren, werden durch oft wiederkehrende Benutzung wichtig. Ein solches Erleichterungsmittel ist eine im Jahr 1812 in der monatlichen Correspondenz
zuerst gegebene Hülfstafel, um aus den Logarithmen zweier Grössen unmittelbar
die Logarithmen ihrer Summe oder Differenz abzuleiten: man erreicht dadurch
mit Einem Aufschlagen, wozu man sonst ein dreimaliges oder wenigstens zweimaliges nöthig hätte. Im Besitz einer solchen Tafel wird man mit Vortheil manche
\SourceLeaf{260}
Formeln in ihrer ursprünglichen Gestalt beibehalten können, denen man sonst
wohl, durch Einführung von Hülfswinkeln, eine zur Rechnung bequemere Form
zu geben sucht. Die erwähnte Tafel war nur für Rechnungen mit fünf Decimalen
bestimmt: eben weil solche Rechnungen am häufigsten vorkommen, lag dies Bedürfniss am nächsten. Der bei Bekanntmachung derselben geäusserte Wunsch,
dass jemand sich der Mühe unterziehen möchte, eine ähnliche Tafel in grösserm
Umfange und mit sieben Decimalen zu berechnen, hat die vorliegenden Tafeln
veranlasst, deren Verfasser, Hr. \textsc{Matthiessen}, sich durch diese mühsame Arbeit
ein Recht auf den Dank aller derer erworben hat, die viel mit logarithmischen
Rechnungen zu thun haben. Die Vorrede gibt Nachricht von der Methode, deren sich der Verfasser zur Berechnung bedient hat, und von seiner lobenswerthen
Sorgfalt, die Tafel auch in der letzten Ziffer durchgehends zuverlässig zu machen.
Bei Tafeln dieser Art ist auch die äussere Einrichtung keinesweges gleichgültig.
Der Verf. hatte anfangs eine Anordnung im Sinn, deren Zweck war, die Tafel in
den möglich kleinsten Raum zusammenzudrängen. Allein da die ganze Bestimmung der Tafel nur dahin geht, die Rechnungen zu \textit{erleichtern}, so würde jene
ganz verfehlt werden, wenn die Anordnung der Tafel zu künstlich wäre, und eine
beschwerliche Aufmerksamkeit erforderte. Der Verf. entschloss sich daher bei
reiferer Ueberlegung mit Recht, jene wenn gleich sinnreiche Anordnung bei Seite
zu setzen, und statt derselben eine einfachere zu wählen, obgleich der Umfang
des Bandes dadurch beträchtlich vergrössert wurde. Vielleicht wird mancher,
der die Tafeln gebraucht, mit uns wünschen, dass der Verf. hierin lieber noch
etwas weiter gegangen wäre, und die Columne \(B\) und \(C\) jede vollständig hätte
abdrucken lassen, deren vier letzte, beiden gemeinschaftliche, Ziffern nur einmal
dastehen. Der Ueberblick würde dadurch noch bequemer geworden sein, und
das Format wäre auch dadurch nicht vergrössert, wenn etwas kleinere Schrift gewählt wäre, welches ohne Nachtheil, vielleicht selbst mit Vortheil für das gefällige Ansehen, hätte geschehen können. Auch die fast zu strenge Oekonomie mit
den Ziffern, wo in der Regel nur die vierte abgedruckt ist, so lange die vorhergehenden ungeändert bleiben, und mit den Proportionaltheilen, die immer nur
Einmal angesetzt sind, und also zuweilen ein Zurück- oder Vorausblättern nöthig
machen, thut der Bequemlichkeit einigen Eintrag. Endlich hätten wir gewünscht,
dass der letzte Theil der Tafel S. 212 in einer zehnmal. grössern Ausdehnung gegeben wäre, um die zweiten Differenzen unmerklich zu machen; es wäre dazu
\SourceLeaf{261}
nur Eine Seite mehr erforderlich gewesen. Alle diese Bemerkungen, die zum
Theil mit auf individueller Gewöhnung beruhen mögen, sollen das Verdienstliche
dieser Arbeit keinesweges schmälern, welches gewiss von allen anerkannt wird,
die von derselben Gebrauch zu machen Gelegenheit nehmen werden.

\par\Needspace{12\baselineskip}\pdfbookmark[1]{Gauss an Schumacher. 1822 November 25.}{schumacher1822}\begin{center}\rule{0.72\linewidth}{0.3pt}\par Astronomische Nachrichten, herausg. v. \textsc{Schumacher}. Nr. 24. Beilage 2. 1822 Dec.\par\rule{0.72\linewidth}{0.3pt}\end{center}

\begin{center}
\textsc{Gauss} an \textsc{Schumacher}.
\end{center}
\begin{flushright}
Göttingen 1822. Nov. 25.
\end{flushright}
--- --- Den von Herrn Professor \textsc{Encke} geäusserten Wunsch Logarithmentafeln mit
6 Ziffern betreffend, habe ich schon öfters gehegt, und Sie werden sich gewiss
vielfältigen Dank erwerben, wenn Sie solche veranlassten. Alles, was Herr \textsc{Encke}
über das Aeussere und Innere sagt, unterschreibe ich als meine eigene Meinung,
nur die Proportionaltheile scheinen mir überflüssig, und alles Ueberflüssige schadet dem leichten übersichtlichen Gebrauch. Bei solchen Dingen hängt freilich
manches von individueller Gewöhnung ab; indessen wenn einige, die Gebrauch
von Tafeln machen, anders gewöhnt sind als ich, so sind doch auch wohl andere
eben so gewöhnt; und daher berühre ich noch einen Umstand, nämlich die Abänderung der \(4^{\mathrm{ten}}\) Ziffer für die Logarithmentafeln. Sie kennen die Einrichtung,
die in dieser Beziehung in \textsc{Callet}'s Tafeln gemacht ist, und einige haben dies als
eine Verbesserung betrachtet. Ich gestehe, dass ich der entgegengesetzten Meinung bin, und die \textit{regelmässige} Abtheilung von 5 zu 5 Zeilen durch horizontale
Striche, wie sie in \textsc{Sherwin}'s und andern Tafeln ist, für etwas, bei häufigem Gebrauche \textit{viel} wesentlicheres und bequemeres halte, daher ich mich der \textsc{Callet}'schen Logarithmen auch niemals bedienen mag. Bei meiner vieljährigen Praxis
weiss ich auch nicht einen einzigen Fall, wo der Gebrauch der \textsc{Sherwin}'schen Tafeln mich bei der \(4^{\mathrm{ten}}\) Ziffer zu einem Rechnungsfehler verleitet hätte, daher ich
auch auf die Sternchen bei \textsc{Vega} und andern Tafeln gar keinen Werth lege, und
des bessern Papiers und der schönern Ziffern wegen, mich lieber an die \textsc{Sherwin}'schen halte. Wer auf solche Warnungszeichen einen Werth setzt, kann sich
leicht in seinem Exemplare an den betreffenden Stellen rothe oder grüne Punkte
machen. ---
\clearpage\fancyhead[L]{\small ANZEIGEN. LOGARITHMISCHE TAFELN.}
\SourceLeaf{262}
\pdfbookmark[1]{Babbage. 1828 Januar 19.}{babbage1828}\par\Needspace{10\baselineskip}\begin{center}\rule{0.7\linewidth}{0.4pt}\par Göttingische gelehrte Anzeigen. 1828 Januar 19.\par\rule{0.7\linewidth}{0.4pt}\end{center}

\noindent\textit{Table of logarithms of the natural numbers, from} 1 \textit{to} 108000, \textit{by} \textsc{Charles}
\textsc{Babbage.} \textit{Stereotyped. London. Printed for J. Mawman.} 1827. 202 S. gr. 8.

Dieser neue Abdruck der Logarithmentafeln zeichnet sich vor andern durch
eine geflissentlichere Beachtung kleiner Nebenumstände aus. Wer nur von Zeit
zu Zeit einmal veranlasst wird, einige Logarithmen in den Tafeln aufzusuchen,
verlangt von ihnen hauptsächlich nur möglich grösste Correctheit. Allein für andere, denen die Tafeln ein tägliches Arbeitsgeräth sind, bleiben auch die geringfügigsten Umstände, die auf die Bequemlichkeit des Gebrauchs Einfluss haben
können, nicht mehr gleichgültig. Farbe, Stärke und Schönheit des Papiers;
Format; Grösse, Schärfe und gefälliger Schnitt der Typen: Beschaffenheit der
Druckerschwärze; Anordnung der Zahlen, um das was man sucht ohne Ermüdung des Auges schnell und sicher zu finden; Vorhandensein von allem, was man
braucht, aber auch Abwesenheit von allem, was man nicht brauchen mag, und
was sonst die leichte Uebersicht nur stören würde, alle diese Umstände erhalten
eine gewisse Wichtigkeit bei einem Geschäfte, welches man täglich hundert mal
wiederholt. Freilich hängt dabei manches von der Individualität des Rechnenden
und von seiner Gewohnheit ab, so dass nicht wohl Eine Ausgabe allen am besten
gefallen kann; dem Kurzsichtigen ist ein grosses Format beschwerlich, und er
zieht kleinere Typen vor, während es sich bei dem Weitsichtigen umgekehrt verhält; der weniger geübte Rechner legt einen Werth auf diesen oder jenen Zusatz,
welchen der geübtere, der mehrere Theile der Operationen ohne Anstrengung im
Kopf macht, lieber wegwünscht, weil alles Ueberflüssige nur störend wirkt. Inzwischen gibt es doch auch allgemeingültige Regeln. Der Herausgeber der vorliegenden Tafeln hat in der Vorrede ein Dutzend solcher Vorschriften zusammengestellt, die er aus der Vergleichung vieler Logarithmentafeln abgeleitet hat, und
die meistens sogleich von selbst einleuchten, obwohl einige davon nur unter Einschränkungen anzuerkennen sein möchten.

Was die gegenwärtige Ausgabe der Logarithmentafeln am meisten von andern unterscheidet, ist, dass sie auf farbiges (gelbes) Papier abgedruckt ist: man
gewöhnt sich daran bald, und findet es, besonders zum Gebrauche bei Licht, an%
\SourceLeaf{263}genehmer als weisses. Die Anordnung ist im Wesentlichen die gewöhnliche.
Nach einem Wechsel der dritten Ziffer ist für die übrigen Logarithmen in derselben Zeile die vierte Ziffer mit andern Typen, etwa halb so gross wie die übrigen,
gesetzt. Für Ref., der überhaupt auf solche Warnungszeichen wenig Werth setzt,
haben diese Typen, auch nach einem Gebrauch von ein paar Monaten, noch nicht
das fremdartig Störende verloren, und er würde den in \textsc{Vega}’s Tafel gebrauchten
Sternchen, oder den Punkten, die in die bessern Exemplare von \textsc{Taylor}’s Tafeln
eingedruckt sind, den Vorzug geben. Sternchen hat der Herausgeber desswegen
nicht gebrauchen wollen, weil sonst die Columnen zu breit geworden, und dann
für die Verwandlung der Zahlen, als Secunden betrachtet, in Grade, Minuten
und Secunden kein Raum geblieben wäre. Referent würde diese Verwandlungscolumne (eben so eingerichtet, wie man sie aus \textsc{Callet}’s Tafeln kennt) gern entbehrt haben; ein geübter Rechner wird nicht, einer so leicht selbst im Kopfe zu
machenden Verwandlung wegen, erst die Tafeln aufblättern; zweckmässig ist es
aber, dass der Herausgeber diese Columnen wenigstens durch eine starke Linie
von der Logarithmentafel geschieden hat. Endlich ist der Fall wo die letzte Ziffer eine Vergrösserung erlitten hat, (weil der weiter fortlaufende Logarithm als
achte Ziffer 5 oder eine grössere gehabt haben würde) durch einen unter diese siebente Ziffer gesetzten Punkt ausgezeichnet; diese Einrichtung ist wenigstens besser, als die von \textsc{Prasse} gewählte, durch Typen von anderer Form; doch behalten
auch so einige Ziffern noch etwas unangenehm Fremdartiges, was nach unserer
schon bei Anzeige der \textsc{Prasse}’schen Tafeln in diesen Blättern [1811 Mai 25] erwähnten Ansicht durch den Nutzen nicht aufgewogen wird.

Uebrigens lassen Typen und Papier bei dieser Ausgabe der Logarithmentafeln nichts zu wünschen übrig, und auf die Correctheit ist die ausgezeichnetste
Sorgfalt verwandt. Bei einer Ausgabe, wo auf die kleinen die Bequemlichkeit
des Gebrauchs angehenden Umstände so viel Sorgfalt verwandt ist, fällt die spärliche Ansetzung der Proportionaltheile auf; häufig muss man, um diejenigen zu
finden, die man eben braucht, voraus- oder zurückblättern.
\SourceLeaf{264}
\pdfbookmark[1]{Hassler. 1831 März 31.}{hassler1831}\par\Needspace{10\baselineskip}\begin{center}\rule{0.7\linewidth}{0.4pt}\par Göttingische gelehrte Anzeigen. 1831 März 31.\par\rule{0.7\linewidth}{0.4pt}\end{center}

\noindent\textit{Tabulae logarithmicae et trigonometricae notis septem decimalibus expressae.}
\textit{In forma minima. Purgatae ab erroribus praecedentium tabularum cura} F. K. \textsc{Hassler.} New York. G. C. und H. Carvill.

Diese transatlantische Ausgabe der Logarithmentafeln zeichnet sich durch
eine ganz vorzügliche Nettigkeit des stereotypisch ausgeführten Drucks aus. Sie
enthält, durchgehends auf sieben Decimalstellen, die Logarithmen der Zahlen
bis 100000, die Logarithmen der Sinus, Tangenten, Cosinus und Cotangenten
im ersten Grade durch alle Secunden, in den beiden folgenden von zehn zu zehn
Secunden, und für alle übrigen Grade des Quadranten von dreissig zu dreissig
Secunden. Ausserdem die natürlichen Sinus und Tangenten durchgehends von
dreissig zu dreissig Secunden. Und diess alles in dem Raum von 312 Seiten in
klein Octav oder gross Duodez. Eine solche Zusammendrängung war freilich nur
durch sehr kleine Typen zu erreichen, welche, bei aller Schönheit, doch wohl
für die meisten nicht kurzsichtigen Augen zum täglichen Gebrauch fast zu klein
sein möchten. Auf die Correctheit scheint eine ganz besondere Sorgfalt gewandt
zu sein, wenigstens ist uns bei dem eine Zeitlang versuchten häufigen Gebrauch
gar kein Druckfehler aufgestossen.

\pdfbookmark[1]{Auflösung quadratischer Gleichungen. 1840.}{quadratisch1840}\par\Needspace{10\baselineskip}\begin{center}\rule{0.7\linewidth}{0.4pt}\par \textsc{Vega} und \textsc{Hülsse.} Sammlung mathematischer Tafeln. 1840.\par\rule{0.7\linewidth}{0.4pt}\end{center}

\noindent\textit{Auflösung quadratischer Gleichungen in der Form, dass nicht die Coëfficienten}
\textit{der Gleichung selbst, sondern deren Logarithmen gegeben sind, und dass man auch}
\textit{nicht ihre Wurzeln selbst (oder eine derselben), sondern vielmehr deren Logarithmen}
\textit{zu anderweitiger Benutzung nöthig hat.}

Die Ausführung dieses Geschäfts, bloss mit Hülfe der gewöhnlichen Logarithmentafeln, erfordert nothwendig ein vierfaches Aufschlagen in denselben; man
reicht mit einem zweifachen Aufschlagen aus, wenn man entweder die trigonometrischen Logarithmentafeln oder die Tafeln für Logarithmen von Summen und
\SourceLeaf{265}
Differenzen benutzt; man reicht mit einem einfachen Aufschlagen aus, wenn letztere Tafel, wie hier, durch die Zusatzcolumnen $D$, $E$ und $F$ erweitert ist.

Die vierte Columne, $D$, enthält $B+C$, die fünfte, $E$, enthält $A+C$,
die sechste, $F$, enthält $B-A$. Nehmen wir an, dass die Zahlen selbst, denen
die Logarithmen $A$, $B$, $C$ zugehören, beziehungsweise $a$, $b$, $c$ sind, so enthielte
also die vierte Columne die Logarithmen von $bc$, die fünfte die von $ac$, die
sechste die von $\frac{b}{a}$.

Es mag noch darauf aufmerksam gemacht werden, dass die Logarithmen in
der sechsten Columne $F$, welche bis zu $A=0{,}208$ positiv sind, von $A=0{,}209$
an negativ werden (S.\ 640); es ist ziemlich gleichgültig, ob sie in dieser negativen Form oder, wie hier, durch ihre Complemente angesetzt werden; also könnte
z.\ B.\ für $A=0{,}367$ unter $F$ stehen $-0{,}21180$ oder, wie hier, $9{,}78820$.

Für die Anwendung dieser Zusatzcolumnen auf die Auflösung der quadratischen Gleichung
\[
pxx+qx+r=0
\]
selbst müssen vier verschiedene Fälle unterschieden werden, nämlich
\begin{enumerate}
\item[I.] $p$ und $r$ haben gleiche Zeichen und $\frac{qq}{pr}$ ist nicht kleiner als $4$;
\item[II.] $p$ und $r$ haben gleiche Zeichen und $\frac{qq}{pr}$ ist kleiner als $4$;
\item[III.] $p$ und $r$ haben entgegengesetzte Zeichen und $-\frac{pr}{qq}$ ist grösser als $2$;
\item[IV.] $p$ und $r$ haben entgegengesetzte Zeichen und $-\frac{pr}{qq}$ ist kleiner als $2$.
\end{enumerate}
Der Fall, wo $-\frac{pr}{qq}=2$ ist, kann sowohl zu III.\ als zu IV.\ gezählt werden.

Im Falle II.\ sind die Wurzeln imaginär, in den übrigen erhält man jede der
beiden Wurzeln auf eine doppelte Art. Durch folgendes Schema ist alles leicht
zu übersehen, wobei
\[
\frac{q}{p}=h,\quad \frac{r}{q}=g
\]
gesetzt ist, theils zur Abkürzung, theils weil die Rechnung wirklich in dieser
Form am bequemsten geführt wird.

\begin{center}
\begingroup\renewcommand{\arraystretch}{1.7}\begin{tabular}{rlll}
 & & \emph{Erste Wurzel.} & \emph{Zweite Wurzel.} \\
I. & $+\frac{h}{g}=bc=d$ & $-\frac{h}{b}=-gc$ & $-gb=-\frac{h}{c}$ \\
II. & $+\frac{h}{g}<4$ & Imaginär. & Imaginär. \\
III. & $-\frac{g}{h}=ac=e$ & $+ha=-\frac{g}{c}$ & $+\frac{g}{a}=-hc$ \\
IV. & $-\frac{g}{h}=\frac{b}{a}=f$ & $+\frac{h}{a}=-\frac{g}{b}$ & $+ga=-hb$
\end{tabular}\endgroup
\end{center}
\par
\SourceLeaf{266}
Die Beweise der Vorschriften wird sich jeder leicht selbst entwickeln können. Ein Beispiel des Gebrauchs mag zum Ueberfluss noch hergesetzt werden.

\begin{center}
\begin{tabular}{rl}
Es sei gegeben & $\log p=0{,}69897$ \\
 & $\log q=0{,}84510$ \\
 & $\log r=0{,}77815\,\mathrm{n}$ \\
Also & $\log h=0{,}14613$ \\
 & $\log g=9{,}93305\,\mathrm{n}$ \\
 & $\log(-\frac{g}{h})=9{,}78692$
\end{tabular}
\end{center}

Man sieht sogleich, dass hier der Fall IV.\ statt findet und also $9{,}78692$ in
der sechsten Columne unter $B-A$ oder $F$ zu suchen sein wird. Was von der
Tafel (S.\ 643) hier nöthig wäre ist nur Folgendes:
\begin{center}
\begin{tabular}{c|c|c}
$A$ & $B$ & $B-A$ oder $F$ \\
$0{,}367$ & $0{,}15520$ & $9{,}78820$ \\
$0{,}368$ & $0{,}15489$ & $9{,}78689$
\end{tabular}
\end{center}
Zu $F=9{,}78692$ gehört also
\begin{center}
\begin{tabular}{l}
$A=0{,}36798=\log a$, woraus $\log\frac{h}{a}=9{,}77815$, $\log ga=0{,}30103\,\mathrm{n}$ oder \\
$B=0{,}15490=\log b$, woraus $\log(-\frac{g}{b})=9{,}77815$, $\log(-hb)=0{,}30103\,\mathrm{n}$.
\end{tabular}
\end{center}

\par\Needspace{12\baselineskip}\pdfbookmark[1]{Vega. Thesaurus Logarithmorum. 1851.}{vega1851}
\begin{center}
\rule{0.7\linewidth}{0.4pt}\par
Astronomische Nachrichten. Nr.\ 756. 1851. Mai 2.\par
\rule{0.7\linewidth}{0.4pt}
\end{center}
\begin{center}\emph{Einige Bemerkungen zu} \textsc{Vega}'s \emph{Thesaurus Logarithmorum, von Herrn}\par
\emph{Geheimen Hofrath} \textsc{Gauss}.\end{center}

Der \emph{Thesaurus Logarithmorum} von \textsc{Vega} ist bekanntlich seinem grössten
Theile nach ein neuer Abdruck der grössern \textsc{Vlacq}'schen Logarithmentafeln. In
der Vorrede führt \textsc{Vega} eine nicht unbeträchtliche Anzahl von Fehlern im Original an, die er verbessert hat, mit dem Zusatz, dass er ausser diesen noch eine
sehr grosse Menge von Unrichtigkeiten an der letzten Ziffer der Logarithmen berichtigt habe, zu dem Betrage von einer, zwei, drei, vier Einheiten. Mit gleicher Sorgfalt seien auch die neu hinzugekommenen Tafeln (namentlich also die in
\SourceLeaf{267}
den beiden ersten Graden für alle einzelnen Secunden angegebenen Logarithmen
der trigonometrischen Grössen) berechnet, geprüft und berichtigt. \textsc{Vega} scheint
nun mit der Hoffnung sich geschmeichelt zu haben, dass auf diese Weise seine
Tafeln fast fehlerfrei geworden seien, und verspricht, um zu \emph{vollkommen} fehlerfreien Tafeln zu gelangen, für die erste Anzeige jedes etwa noch stehen gebliebenen Fehlers, der zu falscher Rechnung Anlass geben könne (\emph{pro sphalmatibus calculum turbantibus}) eine Prämie von einem Ducaten zu bezahlen. Ob diese vom
$1^{\mathrm{sten}}$ October $1794$ datirte Ausgelobung jemals Folge gehabt hat, ist mir nicht
bekannt.

Es ist mir zweifelhaft, ob \textsc{Vega} sich ganz klar gemacht habe, was für Fehler als möglicher Weise zu falschen Rechnungen Anlass gebend betrachtet werden sollten. Für alle Tafeln, welche bestimmt sind, theoretisch feststehende
irrationale Grössen darzustellen, gilt bekanntlich der Grundsatz, dass die Tabulargrösse dem wahren Werthe allemal so nahe kommen soll, als bei der gewählten Anzahl von Decimalstellen möglich ist, und es darf folglich die Abweichung
niemals mehr als eine halbe Einheit der letzten Decimale betragen. Jeder Verstoss gegen diese strenge Norm ist ein Fehler, der möglicherweise einen sich auf
die strenge Uebereinstimmung verlassenden Rechner zu einem unrichtigen Resultate verleiten kann. Lässt man diese strenge Auslegung fahren, und mischt in
sein Urtheil eine Rücksicht auf \emph{Erheblichkeit} der Unrichtigkeit ein, so verirrt
sich die Entscheidung in das Gebiet der Willkühr. Der schon vorhin erwähnte
Umstand, dass \textsc{Vega} selbst von Correctionen an den \textsc{Vlacq}'schen Tafeln spricht,
die nur eine Einheit in der letzten Stelle betrugen, und dass er vollkommene Fehlerfreiheit wie sein Ziel bezeichnet, scheint allerdings darauf hinzudeuten, dass
er die strenge Auslegung im Sinne gehabt. Auch habe ich den ersten Theil des
Thesaurus, der die Logarithmen der laufenden ganzen Zahlen enthält, bei sehr
vielen gelegentlich gemachten Vergleichungen mit mehrstelligen Bestimmungen
immer sehr correct gefunden.

Es sind seit jener Zeit bei mehrern andern logarithmischen Tafeln, in der
Absicht ihre Correctheit zu vervollständigen, ähnliche Ausbietungen von Preisen
für die erste Anzeige von Fehlern in den Zahlen gemacht: ich weiss jedoch nicht,
ob dieselben einen Erfolg gemacht haben, mit Ausnahme des bei Tauchnitz in
Leipzig $1847$ von \textsc{Köhler} herausgegebenen logarithmisch-trigonometrischen Handbuchs. Der Verleger dieser Tafeln versprach bei dem ersten Erscheinen, für die erste
\SourceLeaf{268}
binnen einer gesetzten Frist eingesandte Anzeige eines jeden Fehlers, welcher
falsche Resultate veranlassen könne, einen Louisd'or zu bezahlen, und nach einem gedruckten Bericht vom 1. Juli 1848 ist diese Prämie für vier zur Anzeige
gebrachte Fehler wirklich ausgezahlt. Was nun dabei eine ehrende Erwähnung
verdient, ist der Umstand, dass dem einen dieser Fehler jene Qualification nur
unter Anerkennung obiger strengen Auslegung zugesprochen werden konnte. Es
war nemlich der Logarithm von \(103000\), welcher, auf \(12\) Stellen genau,
\[
=5{,}012837224705
\]
ist, in den ersten Abdrücken mit acht Ziffern
\[
=5{,}01283723
\]
angesetzt, während die principmässig abgekürzte Zahl \(5{,}01283722\) ist. Es ist
zu wünschen, dass in künftig vorkommenden ähnlichen Fällen, diese Entscheidung als Präcedenz respectirt werde.

Bekannte sich \textsc{Vega} zu derselben strengen Auslegung. so hätte es ihm leicht
gehen mögen, wie dem König Shiram, dessen Kornkammern nicht ausreichten,
dem Erfinder des Schachspiels die ihm zugesagte Belohnung zu gewähren.

Dass die von \textsc{Vega} ausgebotene Belohnung, unter jener Voraussetzung. ihm
theuer zu stehen kommen konnte, lässt sich schon, ohne alles Nachrechnen, aus
einem Umstande erkennen, der leicht zu bestätigen, jedoch meines Wissens anderweit noch nirgends zur Sprache gebracht ist. Dieser Umstand besteht darin,
dass in der Tafel für die Logarithmen der trigonometrischen Grössen die Zahlen
der Sinuscolumne fast ohne Ausnahme\HistoricalNote{21} der Summe der Zahlen der Cosinuscolumne und der Tangentencolumne genau gleich sind. Da alle diese Zahlen nur
abgekürzte Werthe der irrationalen genauen Grössen sind, so ist klar, dass bei
streng richtiger Abkürzung jene Gleichheit nicht Statt finden wird, in allen den
Fällen, wo die Abweichungen von den genauen Werthen in der zweiten und dritten Columne gleiche Zeichen haben, und ihre Summe mehr als eine halbe Einheit der letzten Decimale beträgt. Man übersieht leicht, dass bei einer grossen
\SourceLeaf{269}
Menge von Fällen dieser Ausnahmefall durchschnittlich einmal unter vieren vorkommen wird, was sich bei wirklicher Abzählung in solchen siebenziffrigen Tafeln, wo auf die Richtigkeit der letzten Ziffer mit Sorgfalt gehalten ist, bestätigt
findet. (Es hat z. B. eine solche Abzählung in den erwähnten \textsc{Köhler}'schen Tafeln an den \(900\) einzelnen Minuten von \(30^{0}0'\) bis \(44^{0}59'\) genau \(225\) solcher Fälle
ergeben, welche vollkommene Uebereinstimmung allerdings für zufällig zu halten ist). \textsc{Vega}'s Thesaurus enthält die Logarithmen der Sinus, Cosinus und Tangenten von \(22680\) Bogen. Unter der Voraussetzung also, dass alle Zahlen in
zweien dieser Columnen scharf nach dem Princip abgekürzt seien, wird man, mit
einer geringen Unsicherheit im Mehr oder Weniger, in der dritten Columne \(5670\)
Logarithmen erwarten dürfen, die um eine Einheit unrichtig angesetzt sind. Für
den Preis von so vielen Ducaten würden wohl befähigte Personen zur vollständigen Neurechnung bereit gewesen sein.

Welche zwei Columnen die ursprünglichen sind, wird man aus der Vergleichung mit anderweitigen auf mehr als zehn Stellen zuverlässigen Bestimmungen erkennen können, schon an wenigen Bögen und mit Gewissheit, wenn die
Voraussetzung der strengen Richtigkeit der ursprünglichen Columnen zutrifft, im
entgegengesetzten Fall aber an einer etwas beträchtlichern Menge wenigstens mit
überwiegender Wahrscheinlichkeit.

Einiges hiezu dienliche findet man schon fertig vor am Schluss der Decimaltafeln, welche \textsc{Hobert} und \textsc{Ideler} 1799 geliefert haben, und die zwar nur mit
sieben Ziffern abgedruckt sind, aber in der Handschrift auf doppelt so viele vollendet waren. Es werden daselbst \(138\) fehlerhafte Logarithmen des \textsc{Vega}'schen
Thesaurus angezeigt, von denen \(127\), um genau gesetzmässig zu werden, einer
Correction von einer Einheit in der zehnten Stelle bedürfen, \(10\) einer Correction von
zwei Einheiten, und eine drei Einheiten. Von den \(127\) fallen \(49\) auf Sinus, \(35\)
auf Cosinus und \(43\) auf Tangenten; die übrigen \(11\) Correctionen, von \(2\) oder \(3\)
Einheiten, beziehen sich blos auf Tangenten. Die betreffenden Bögen sind die
mit \(27\) Minuten messbaren, und die nicht mit angeführten sind diejenigen, wo
die Logarithmen keiner Correctionen bedurften. Es hätten übrigens diese Vergleichungen auch ohne die \textsc{Hobert-Ideler}'schen handschriftlichen Tafeln gemacht
werden können, da die allgemein verbreiteten \textsc{Callet}'schen Tafeln die Logarithmen der Sinus und Cosinus auf \(14\) Ziffern für alle Tausendtheile des Quadranten
enthalten. Wenn man aus diesen auch noch das Nöthige für alle Bögen unter
\SourceLeaf{270}
zwei Grad entlehnt, die durch \(5'24''\) aber nicht durch \(27'\) messbar sind, so gewinnt man die Vergleichung \textsc{Vega}'scher Logarithmen von \(118\) Bögen mit schärfern Bestimmungen, wovon die Resultate in folgendem Abriss (I) zusammengestellt sind:
\[
\begin{array}{rccc}
 & \mathrm{Sin.} & \mathrm{Cos.} & \mathrm{Tang.} \\
0 & 65 & 75 & 54 \\
1 & 53 & 43 & 53 \\
2 &    &    & 10 \\
3 &    &    & 1
\end{array}
\]
Die Bedeutung dieser Zahlen, z. B. der in der letzten Columne ist, dass unter
\(118\) Tangenten-Logarithmen \(54\) richtig angesetzt sind, \(53\) einer Correction von
einer Einheit in der zehnten Stelle bedürfen, \(10\) einer Correction von zwei Einheiten, und einer der Correction von drei Einheiten. Es folgt hieraus schon entschieden, dass \emph{keine} Columne der \textsc{Vega}'schen Tafel durchaus richtig angesetzt
ist, und mit überwiegender Wahrscheinlichkeit, dass die Tangenten-Logarithmen
die abgeleiteten sind, durch die einfache Subtraction der Zahlen der Cosinuscolumne von denen der Sinuscolumne.

Die Verfasser der erwähnten Decimaltafeln hätten übrigens doppelt so viele
Vergleichungen aus ihrer Handschrift geben können. Man kann sich aber eine
noch viel grössere Ausbeute verschaffen, wenn man die höchst schätzbare Tafel
von \textsc{Brigg}'s (in der nach dessen Tode von \textsc{Gellibrand} 1633 herausgegebenen Trigonometria Britannica) benutzt, welche die Logarithmen der Sinus und Cosinus
für alle Hunderttheile der gewöhnlichen Grade auf \(14\) Ziffern liefert, und aus
welcher \textsc{Callet} die oben erwähnten Zahlen entlehnt hat. Sie enthält das Material, um \textsc{Vega}'s Tafeln bei \(1060\) Bögen, also zusammen \(3180\) Logarithmen, prüfen zu können. Ich selbst habe mich jedoch darauf beschränkt, diese Prüfung
an \(81\) Bögen oder an \(243\) Logarithmen, von \(14^{0}0'\) bis \(18^{0}0'\) vorzunehmen, wovon
das Resultat hier folgt (II):
\[
\begin{array}{rccc}
 & \mathrm{Sin.} & \mathrm{Cos.} & \mathrm{Tang.} \\
0 & 29 & 56 & 36 \\
1 & 52 & 25 & 42 \\
2 &    &    & 3
\end{array}
\]
\par
\SourceLeaf{271}
Wirft man die Gruppen (I) und (II) zusammen, so jedoch, dass man die
beiden gemeinschaftlichen Bögen nur einmal in Rechnung bringt, so erhält man
für 190 Bögen folgende Ausbeute (III):
\begin{center}
\begin{tabular}{rrrr}
 & Sin. & Cos. & Tang. \\
0 & 89 & 126 & 87 \\
1 & 101 & 64 & 90 \\
2 & & & 12 \\
3 & & & 1
\end{tabular}
\end{center}
Will man diese Zahlen zu einer Abschätzung des Verhältnisses zwischen den
richtigen und unrichtigen Logarithmen anwenden, so muss man erst noch einen
kleinen Abzug machen. Es ist in (I) der Logarithm des Sinus von $45^{0}$ zweimal
gezählt, nemlich zugleich auch als $\log\cos45^{0}$; es ist ferner in der Tangentencolumne auch der Logarithm der Tangente von $45^{0}$ mitgezählt, der doch rational
ist. Dasselbe gilt von (III) und man muss daher, zu obigem Zweck, das Resultat davon so aussprechen, dass unter 568 geprüften irrationalen Logarithmen 301
sich als richtig, und 267 als unrichtig ausgewiesen haben. Dürfte man dies Verhältniss als durchschnittlich zutreffend betrachten, so würden unter den 68038
irrationalen Logarithmen des \textsc{Vega}'schen Thesaurus (indem man, wie billig, die
Cotangenten nicht mitzählt) nach der Wahrscheinlichkeit etwa 31983 fehlerhafte
anzunehmen sein.

Wahrscheinlich ist aber diese Zahl noch bedeutend zu klein. Ich finde in
meinen Papieren die vor längerer Zeit und zu andern Zwecken auf 14 Ziffern gemachte Berechnung der trigonometrischen Logarithmen für ein paar Gruppen von
Bögen, die nicht sprungsweise, sondern in denselben Intervallen wie die \textsc{Vega}'schen Tafeln, fortschreiten, woraus wenigstens hervorgeht, dass obiges Resultat
(III) noch keinen richtigen Maassstab für die Ungenauigkeit dieser Tafeln abgibt.
Während bei jenen 190 Bögen kein einziger Fall vorkommt, wo der Logarithm
eines Sinus oder eines Cosinus um mehr als \textit{eine} Einheit verbessert werden müsste,
sind solche Fälle gar nicht selten bei denjenigen Bögen, die \textit{nicht} in der Trigonometria Britannica vorkommen, und wo also das zur Vergleichung nöthige erst
durch besondere Rechnung herbeigeschafft werden muss. Ich füge daher die Resultate jener Rechnungen hier bei, da sie dazu dienen können, unserer Vorstellung von dem Grade der Ungenauigkeit der Zahlen in \textsc{Vega}'s Thesaurus eine festere Haltung zu geben.
\par
\SourceLeaf{272}
Die Vergleichung der Zahlen im Thesaurus bei den 21 Bögen von $15^{0}38'20''$ bis $15^{0}41'40''$ mit den schärfer berechneten hat ergeben (IV):
\begin{center}\begin{tabular}{r r r r}
& Sin. & Cos. & Tang.\\
0&4&12&1\\
1&9&8&8\\
2&6&1&6\\
3&2&&4\\
4&&&2
\end{tabular}\end{center}
Die Fälle, wo die grössten Abweichungen vorkommen, sind:
\[
\begin{array}{l l r r r}
15^{0}40'20'',&\text{wo die Correctionen}&+3,&-1,&+4\\
15\quad41\quad30,&\text{wo folgende}&+3,&0,&+4
\end{array}
\]
an die Logarithmen des Sinus, des Cosinus und der Tangente angebracht werden müssen.

\Needspace{11\baselineskip}Ebenso hat die Vergleichung der \textsc{Vega}'schen Zahlen bei den 93 Bögen, von $1^{0}19'52''$ bis $1^{0}21'24''$ mit schärferer Rechnung folgende ergeben (V):
\begin{center}\begin{tabular}{r r r r}
& Sin. & Cos. & Tang.\\
0&38&30&17\\
1&39&56&41\\
2&15&7&22\\
3&1&&11\\
4&&&2
\end{tabular}\end{center}
Die grössten Abweichungen finden Statt bei den Bögen $1^{0}20'10''$ und $1^{0}20'15''$; und die nöthigen Correctionen betragen bei ersterm $-2,+1,-4$, bei dem andern $-2,+2,-4$.

Die Gruppen IV und V haben nicht gleichen Ursprung, da die Zahlen des Thesaurus zu der Gruppe IV aus \textsc{Vlacq}'s Trigonometria Artificialis genommen sind, die andern hingegen zu den neu hinzugekommenen gehören, welche unter \textsc{Vega}'s Leitung und Aufsicht von dem Lieutenant \textsc{Dorfmund} berechnet sind. Nach obigen Abrissen erscheinen die letztern wie etwas weniger ungenau als die erstern, wiewohl die Gruppe IV zu wenig zahlreich ist, um ein sicheres Urtheil zu begründen. Jedenfalls folgt aus dem Zusammenwerfen beider Gruppen eine Schätzung
\SourceLeaf{273}
für die Totalungenauigkeit der Tafeln, die eher etwas zu günstig sein wird, als umgekehrt. Aus dieser Vereinigung folgt, für 114 Bögen (VI):
\begin{center}\begin{tabular}{r r r r}
& Sin. & Cos. & Tang.\\
0&42&42&18\\
1&48&64&49\\
2&21&8&28\\
3&3&&15\\
4&&&4
\end{tabular}\end{center}
Es sind also hier unter 342 Logarithmen nur 102. die keiner Verbesserung bedürfen, gegen 240 ungenaue. Nach diesem Verhältniss würde man unter den 68038 irrationalen Logarithmen 47746 ungenaue erwarten können.

Die Summe der Quadrate der Abweichungen findet sich für die Sinus 159, für die Cosinus 96, für die Tangenten 360. Als mittlern Fehler mag man also annehmen für die Sinus $1{,}18$, für die Cosinus $0{,}92$, für die Tangenten $1{,}78$. Man kann hienach nicht zweifeln, dass die Tangenten-Logarithmen die abgeleiteten sind.

\begin{samepage}
Dass die Zahlen der Cosinuscolumne weniger ungenau sind, als die der Sinuscolumne, rührt wohl ohne Zweifel wenigstens theilweise, daher, dass bei den erstern die zur Ausfüllung erforderlichen Interpolationsmethoden einfacher ausfallen, möglicherweise können indess noch andere Ursachen mitgewirkt haben, worüber sich nur unsichere Vermuthungen aufstellen lassen würden.
\begin{center}\bigskip\rule{.3\linewidth}{.4pt}\end{center}
\end{samepage}
